\section{Vignettes del futuro y Q\&A: de la inteligencia barata a la inteligencia social}

Las últimas cuatro diapositivas de la presentación formal extrapolan al futuro el diseño de sistemas expuesto antes: la caída del costo de las capacidades, las interfaces generadas dinámicamente, la personalización de la educación y la salud, y el cambio en la relación con la creación de historias. Después, los 22 minutos de Q\&A aportan condiciones restrictivas más valiosas. Las notas leen juntas la visión y las respuestas: a cada «se puede» le agregan «qué evidence falta».

\subsection{El costo de las capacidades baja, pero el costo de verificación no necesariamente baja al mismo ritmo}

La gráfica de la MMLU cost frontier expresa una tendencia histórica: con una precisión similar, el precio de la inferencia baja rápidamente. No demuestra que el costo vaya a bajar exponencialmente para siempre, ni que MMLU represente el trabajo real. Más importante aún, que generar sea más barato aumenta el volumen de salidas; si la verificación sigue dependiendo de expertos, el costo total del sistema puede desplazarse de la inference a la review.

\lecturefigure{slide-57-mmlu-cost-frontier.jpg}{MMLU performance--cost frontier: caída histórica del costo de invocación para capacidades equivalentes}{00:45:46--00:47:40}

Sea $q(c)$ la tasa de éxito de la tarea, $c$ el costo de una sola inferencia y $v$ el costo de la verificación humana; el costo por éxito confiable puede escribirse, de forma aproximada, como
\[
C_{\text{trusted}}=\frac{c+v}{q(c)}.
\]
Aunque $c$ baje, si las tareas más complejas hacen subir $v$ o vuelven inestable $q$, los resultados confiables no se abaratan al mismo ritmo. El producto debe optimizar $C_{\text{trusted}}$, y no solo reportar el token price.

\teachervoice{00:45:46--00:47:40, la ponente dice que esta gráfica ya es «de hace un año» y que MMLU ya no es el benchmark que más interesa; la usa para expresar que la raw intelligence se abarata. Las notas advierten: es una dated trend, no una ley de extrapolación ilimitada.}

\subsection{Dynamic UI, servicios personalizados y la relación con las historias}

La sección anterior mostró que baja el costo de invocar la raw intelligence, pero una salida barata solo genera valor para el usuario si llega a una interfaz adecuada. Esta sección conecta la tendencia de costos con tres tipos de surface: la UI dinámica generada según la intención, los puntos de acceso a la educación y la salud que asumen recomendaciones de alto riesgo, y la colaboración narrativa que preserva la agency del autor. La visión de la Dynamic generative UI es que la interfaz se genere al instante según la intención: quien aprende de forma visual ve gráficos manipulables; quien aprende de forma auditiva recibe audio. Esto exige que el modelo infiera a la vez la tarea, las preferencias del usuario y los dispositivos disponibles; una inferencia errónea debe poder deshacerse con facilidad, pues de lo contrario el «software invisible» también volverá impredecible el comportamiento del sistema.

\lecturefigure{slide-58-dynamic-generative-ui.jpg}{Dynamic generative UI: la interfaz se genera según la intención y la creación de software tiende a volverse invisible}{00:47:40--00:48:04}

\lecturefigure{slide-59-personalized-health-education.jpg}{Salud y educación personalizadas con consumer hardware como punto de acceso futuro}{00:48:04--00:48:42}

La salud y la educación tienen un gran valor, y también riesgos distintos. Un sistema educativo puede admitir la exploración y la corrección de errores; las recomendaciones médicas implican omisión de síntomas, triaje de urgencias, privacidad y regulación. Un indicador unificado de «personalización» oculta los niveles de riesgo. El producto debe distinguir entre information, decision support y autonomous action, y prever revisión profesional y escalamiento de emergencia para las salidas de alto riesgo.

\lecturefigure{slide-60-storytelling-relationship.jpg}{La relación de las personas con el storytelling process cambiará con los modelos colaborativos}{00:48:42--00:49:18}

La apertura de la creación narrativa la convierte en un buen entorno para poner a prueba la social intelligence, la voice preservation y la agency. El modelo puede dar inspiración, reescrituras y sugerencias de estructura, pero «decidir por el autor qué vale la pena expresar» y «ayudar al autor a realizar su intención» son productos distintos. Un verdadero collaborator debe permitir que el autor vea las fuentes, las opciones y los cambios, en lugar de arrebatarle el control creativo con texto fluido.

\begin{warningbox}{Una future vignette no equivale a una deployment claim}
Estas cuatro diapositivas expresan el juicio de la ponente sobre la dirección a seguir. No demuestran que ya estén resueltas la seguridad médica, la ganancia educativa, la previsibilidad de la UI dinámica ni el impacto sobre el trabajo creativo. Las notas las tratan como una research agenda: cada visión debe completarse con usuarios objetivo, modos de falla, permisos, verificación, costo y responsables.
\end{warningbox}

\teachervoice{00:47:40--00:49:18, Karina imagina interfaces dinámicas guiadas por preferencias visuales o auditivas, educación y salud personalizadas y una nueva relación con el storytelling, y a la vez espera que los creadores vean la IA como una herramienta y no solo como una amenaza. Nota de clase: esta visión optimista debe realizarse junto con la agency del usuario y los mecanismos de verificación.}

\subsection{Q\&A I: tareas subjetivas, benchmarks y rutas de emprendimiento}

El Q\&A pregunta primero: ¿hay tareas difíciles de evaluar que, por ello, los investigadores prefieren no abordar? La ponente pone como ejemplos la creative writing y la emotional intelligence; considera que es posible construir benchmarks, pero las tareas complejas avanzan más despacio y dependen más de hitos. También señala que quien emprende no necesita entrenar primero un foundation model: puede usar modelos existentes de bajo costo y concentrar la innovación en tareas reales y en el ciclo cerrado del producto.

\begin{knowledgebox}{Nota de clase: subjetivo no significa imposible de evaluar}
00:49:18--00:55:00, Karina distingue entre «no existe un frontier benchmark disponible» y «en principio no se puede construir una evaluación». Se puede partir de comparaciones por pares, panels de expertos, outcomes de largo plazo de los usuarios, resultados de competencias y ejemplos contrafactuales; pero todas estas señales conllevan juicios de valor y no deben hacerse pasar por una verdad natural.
\end{knowledgebox}

\subsection{Q\&A II: cómo evitar que las preferences converjan al promedio global}

La ponente propone dos líneas: preservar la entropy del base model y usar RLAIF/synthetic preference data para generar de forma dirigida la diversidad necesaria. Aquí la entropy no es un estímulo a la aleatoriedad, sino una forma de evitar que el RL comprima muchos estilos útiles en «los emojis, el Markdown y el tono que le gustan al usuario promedio».

Para un contexto $x$, la entropy de la política es
\[
\mathcal{H}(\pi(\cdot\mid x))=-\sum_y \pi(y\mid x)\log \pi(y\mid x).
\]
Una preference optimization demasiado intensa puede reducir $\mathcal{H}$ y provocar mode collapse. En la práctica, conviene mantener la diversidad de candidatos según la tarea y, a la vez, filtrar con restricciones los errores factuales y los riesgos de seguridad; «diverso» no significa «cualquier salida es buena».

\begin{importantbox}{Nota de clase: el valor de RLAIF es la cobertura controlable, no un volumen ilimitado de datos}
00:55:00--01:01:20, Karina dice que los synthetic data no necesariamente tienen que ser muchos; lo decisivo es la diversity. El equipo todavía puede revisarlos manualmente, o validarlos con un modelo independiente bien calibrado y con meta-eval. Describe la synthetic generation como una curation de la distribución objetivo, no como una copia de la preferencia humana promedio.
\end{importantbox}

\subsection{Q\&A III: diagnóstico cualitativo, costos y frontier work}

Sobre los bugs del modelo, la ponente reconoce que mucha nuanced weirdness surge de que las personas «juegan con el modelo» una y otra vez; las eval automáticas sirven para revisar behaviors conocidos, pero los problemas nuevos suelen descubrirse primero mediante exploración cualitativa. También distingue el costo del desarrollo de producto del de la frontier research: los equipos de aplicación pueden reutilizar modelos baratos; los equipos de frontera deben asumir la ineficiencia de la etapa de invención, y solo después llega una segunda ronda de optimización de costos.

\begin{knowledgebox}{Nota de clase: una anomalía aislada frente a un comportamiento sistemático}
00:57:20--01:00:00, Karina dice que la revisión cualitativa debe fijarse en si el behavior aparece de forma persistente. Un sample ocasional puede deberse a la aleatoriedad; solo cuando se reproduce de forma estable entre variantes de prompt, sesiones y versiones debe escalarse a issue sistemático. Las notas recomiendan ampliar cada hallazgo manual a un conjunto de variantes y luego incorporarlo a la regresión automática.
\end{knowledgebox}

\begin{warningbox}{Límites de la discusión sobre costos}
Las respuestas del Q\&A sobre el costo de entrenamiento, el scaling y el abaratamiento futuro tienen una incertidumbre evidente. Lo que puede confirmarse es que la capa de aplicación puede reutilizar modelos existentes, que la invención de frontera es costosa y que la ingeniería posterior reduce costos; de las respuestas en clase no puede deducirse que el costo de entrenamiento vaya a bajar necesariamente según una curva determinada.
\end{warningbox}

\subsection{Q\&A IV: robots, AI coworker y social intelligence}

La ponente es optimista respecto de la robotics, pero considera que los datos son un enorme bottleneck. A la pregunta de si «ya existen colegas de IA», responde con claridad que todavía no; la forma ideal se parecería más a compartir pantalla, conversar en tiempo real, señalar objetos concretos y ajustar la agency según la voluntad del usuario. Lo que falta no es otro Slack bot, sino social intelligence, speech-to-speech, grounding multimodal y un espacio de trabajo compartido.

\begin{knowledgebox}{Nota de clase: qué le falta todavía al AI coworker}
01:05:00--01:10:00, las carencias que enumera Karina incluyen: generación en tiempo real y edición conjunta, saber cuándo orientar sin quitar el control, speech-to-speech, señalar al mismo tiempo el objeto del que se habla y adaptarse a distintas relaciones. Esta lista es más exigente que «poder hacer tareas automáticamente» y explica también por qué una teammate eval debe incluir la coordinación interpersonal.
\end{knowledgebox}

\subsection{Q\&A V: Research-driven product frente al flujo tradicional de PRD}

El desarrollo de producto tradicional suele empezar con el PRD, el diseño y la planeación de ingeniería; un research-driven product puede surgir primero de una capability demo sorprendente y luego dar forma al producto en torno a ella, o bien producto e investigación pueden experimentar juntos desde el primer día. La ponente presenta Canvas como ejemplo de lo segundo: el post-training trabaja junto con producto y el proceso es más ad hoc, pero eso no significa que no haya normas; las normas se trasladan a los prototypes, las eval, las traces y la evidencia de iteración.

\begin{importantbox}{Nota de clase: una demo de capacidad no es el punto final del producto}
01:09:52--01:11:29, Karina contrasta el flujo de software tradicional con el research-driven product: cuando la investigación produce primero una capacidad, el producto se forma en torno a ella; a veces producto e investigación se diseñan juntos desde el inicio. El criterio de ingeniería es que, cuanto más ad hoc sea el proceso, más necesario es registrar la hypothesis, las versiones y los failures, para evitar que el equipo solo recuerde las demos exitosas.
\end{importantbox}

\subsection{Resumen de la sección}

La condición común de las visiones de futuro no es «modelos más inteligentes», sino que baje el costo por éxito confiable, que las interfaces puedan verificarse y deshacerse, que la personalización respete el control del usuario, que las tareas subjetivas conserven la diversidad y que el AI coworker tenga capacidades sociales y multimodales. El Q\&A completa la tesis del curso: muchas capacidades pueden enseñarse, pero los datos, la verificación, la distribución de preferencias y la relación de colaboración determinan si se convierten en un producto confiable.
